% Generated draft: requires mathematical and source review.
\documentclass[11pt]{article}
\usepackage[margin=25mm]{geometry}
\usepackage{fontspec}
\setmainfont{Latin Modern Roman}
\usepackage{amsmath,amssymb,mathtools,unicode-math}
\setmathfont{Latin Modern Math}
\usepackage{enumitem,needspace,xurl,hyperref}
\setlength{\parindent}{0pt}
\setlength{\parskip}{5pt}
\setlength{\emergencystretch}{3em}
\newcommand{\N}{\mathbb{N}}
\newcommand{\Z}{\mathbb{Z}}
\newcommand{\Q}{\mathbb{Q}}
\newcommand{\R}{\mathbb{R}}
\newcommand{\C}{\mathbb{C}}
\newcommand{\F}{\mathbb{F}}
\newcommand{\A}{\mathbb{A}}
\newcommand{\PP}{\mathbb{P}}
\newcommand{\E}{\mathbb{E}}
\newcommand{\eps}{\varepsilon}
\newcommand{\dd}{\,\mathrm{d}}
\begin{document}
\section*{UnsolvedMath ID: 30006021}
\textbf{Source problem code:} catalogue entry
\subsection{Definitions and mathematical statement}
A square-tiled surface is a surface assembled from finitely many unit Euclidean squares by identifying boundary edges according to the gluing conventions of the chosen model. The supplied record does not specify these conventions or the probability distribution on such surfaces.

Let \(g\in\N\) denote genus and \(N\in\N\) the number of squares. Fix \(\alpha>0\) and put
\[
N_g=\lfloor\alpha g\rfloor.
\]
The requested problem is to determine the asymptotic behavior of random square-tiled surfaces of genus \(g\) with exactly \(N_g\) squares as \(g\to\infty\), for genera for which the intended ensemble exists. This is a simultaneous limit, in contrast to the order of limits attributed to Theorem 2 in the source: first \(N\to\infty\) at fixed genus, followed by \(g\to\infty\).

More formally, if \(\mathcal S_{g,N}\) denotes the intended ensemble and \(\mu_{g,N}\) its intended probability measure, the question concerns the asymptotic behavior of
\[
X_g\sim\mu_{g,N_g},\qquad X_g\in\mathcal S_{g,N_g}.
\]
Here \(\mathcal S_{g,N}\) and \(\mu_{g,N}\) are placeholders for missing source specifications, not definitions of a particular random model. The record supplies neither the observables and their normalizations nor a topology or other precise meaning of limiting behavior. It also does not specify restrictions on \(\alpha\) ensuring nonempty ensembles. Consequently, this is a qualified research question rather than a fully specified convergence assertion.

The original excerpt also includes a separately numbered Problem 2 concerning cylinders of intermediate area. It attributes an understanding of macroscopic cylinder areas, proportional to the total surface area, and microscopic cylinder areas, proportional to total area divided by \(g\), to Theorem 2. Its proposed intermediate range contains corrupted square-root notation, apparently suggesting normalized areas between \(1/\sqrt g\) and \(2/\sqrt g\); this reading is uncertain. That neighboring question supplies context but does not determine the observables intended in the diagonal-limit problem.
\subsection{Short English statement}
For fixed \(\alpha>0\), study random square-tiled surfaces as their genus \(g\) and number of squares \(\lfloor\alpha g\rfloor\) grow together. A precise limiting statement requires the missing random model, admissibility conditions, observables, and normalizations. The supplied record labels the problem open; that status has not been independently verified.
\subsection{Sources}
\begin{itemize}
\item Local UnsolvedMath registry: supplied complete record with permanent ID \(30006021\), title \emph{Diagonal Large-Genus Limits of Square-Tiled Surfaces}, exact source problem code \texttt{catalogue entry}. The excerpt identifies the diagonal-limit question as Problem 3 and also contains an adjacent Problem 2. No individual proposer is supplied; the proposed year is recorded as 2024.
\item Supplied report citation: \emph{Riemann Surfaces: Random, Flat, and Hyperbolic Geometry} (2025), Oberwolfach Reports, pp. 2369--2424, DOI \texttt{10.4171/owr/2024/41}. Source URL and supplied literature-source link labeled ``OWR report'': \url{https://doi.org/10.4171/owr/2024/41}. The supplied background describes Oberwolfach Reports as documenting mathematical workshops and their open-problem discussions. The report and its referenced Theorem 2 have not been independently consulted.
\item Supplied literature assessment, checked exactly \texttt{2026-08-22}: the diagonal regime is described as explicitly proposed future work, with no primary result located in that curation establishing its limiting behavior. The supplied status is ``open.'' This is supplied dated curation, not an independently verified literature assessment; time-sensitive status and cited claims require independent verification.
\item The registry marks the statement \texttt{corrected\_verified} and states that only formula delimiters, notation, grammar, and typography were repaired while preserving mathematical content and scope. This supplied editorial label does not establish mathematical completeness or constitute independent mathematical review.
\end{itemize}
\end{document}
